\newpage\handout
{Part 1: Enumeration}
{\textsc{Combinatorics (COMB.)}}
{\href{https://creativecommons.org/licenses/by-nc-sa/4.0/}{CC BY-NC-SA 4.0 license}}
{Author: \BookAuthor}{Release: \generatedOn}
{SAT: Part-1}
{\small\textbf{Purpose and Usage:}}\\[0.2cm]
{\footnotesize
These notes represent individual secondary work created by a TA at 
\WorksheetInstitution \space
based on experience with institutional course materials. While derived 
from \WorksheetInstitution \space
content, this compilation includes original explanations, 
examples, and pedagogical approaches developed independently. This work 
is shared as open source educational material and is not officially 
endorsed or authorized by \WorksheetInstitution.}\\[0.5cm]

{\small\textbf{For Educational Use:}}\\[0.2cm]
{\footnotesize
Students and educators are welcome to use and adapt these materials for 
non-commercial educational purposes. Please acknowledge both the original 
\WorksheetInstitution \space
source materials and the individual contributions of the 
author. This work is subject to removal or modification if \WorksheetInstitution 
determines it violates their Honor Code or intellectual property policies.}\\[0.5cm]

{\small\textbf{Content and Attribution:}}\\[0.2cm]
{\footnotesize
This work combines concepts and approaches learned through \WorksheetInstitution 
teaching experience with original pedagogical content. While inspired by 
institutional materials, all explanations, examples, and exercises are 
independently created. This compilation is shared under Creative Commons 
Attribution-NonCommercial-ShareAlike (CC BY-NC-SA) license, with proper 
attribution to both \WorksheetInstitution \space
and the individual author.}\\[0.5cm]

{\small\textbf{Quality Assurance:}}\\[0.2cm]
{\footnotesize
Efforts have been made to ensure clarity and correctness. Users are 
encouraged to verify concepts independently and report any errors or 
concerns. This work is provided as-is for educational benefit and 
may be subject to removal if institutional policies are violated.}

\begin{center}
  \includegraphics[width=0.30\textwidth]{\WorksheetEmblem}\\
  {\small © \the\year\ \BookAuthor. All rights reserved.}\\[0.5cm]
  {\small Prepared by \BookAuthor, TA@GT}\\[1cm]
\end{center}

\newpage

\begin{enumerate}
  \item \textbf{Birthday Triplet Guarantee} (10 points)\\
  A researcher asked his assistant to provide a list of babies born in July for his research project. He is trying to see if three babies share a birthday. What is the minimum number of babies on the list provided by his assistant to guarantee that the researcher can find three babies who share the same birthday?
  \begin{subanswer}
    % your answer here
  \end{subanswer}

  \item \textbf{Candy Purchase Ratio} (10 points)\\
  Sam is buying two kinds of candies: jelly beans and lemon drops. Each jelly bean costs 35 cents, and each lemon drop costs 60 cents. The ratio of the amount he paid for jelly beans to lemon drops is $\frac{7}{8}$. If he paid more than $\$10$, what is the least possible amount he paid in cents? (Ignore taxes)
  \begin{subanswer}
    % your answer here
  \end{subanswer}

  \item \textbf{Rational Numbers with Product 2024} (10 points)\\
  $\frac{a}{b}$ is a rational number written in lowest terms, where $a$ and $b$ are positive integers. If $ab=2024$, how many of these rational numbers are possible?
  \begin{subanswer}
    % your answer here
  \end{subanswer}

  \item \textbf{Sum of Consecutive Numbers} (10 points)\\
  2024 can be written as the sum of $n$ consecutive natural numbers, where $n$ is greater than 1. What is the sum of all possible values of $n$?
  \begin{subanswer}
    % your answer here
  \end{subanswer}

  \item \textbf{Pairwise Remainder Reduction} (10 points)\\
  Natural numbers from 1 to 2024 are written in sequence. First, for each pair of consecutive numbers from left to right, we replace the pair with the remainder of its sum when divided by 9. For example, the pairs $(1,2),(3,4),(5,6)$, and $(7,8)$ are replaced by $3,7,2$, and $6$, respectively.

  Next, we make pairs from right to left and do the same for each pair. Then, we make pairs from left to right and do the same for each pair. In each step, if the last number is not paired with any number, the number is left unchanged. If we continue these steps, what is the last number remaining?
  \begin{subanswer}
    % your answer here
  \end{subanswer}

  \item \textbf{Base Representation Ending Digits} (10 points)\\
  Let $N$ be a positive integer. $N$ ends with 101 when written in base 2 representation. It also ends with 243 when written in base 5 representation. Find the last three digits when $N$ is written in base 10 representation.
  \begin{subanswer}
    % your answer here
  \end{subanswer}

  \item \textbf{Car Grouping at Airport} (10 points)\\
  Six people arrived at Honolulu airport and rented three identical cars. Because of their luggage, each car can take no more than 3 people. If each person can drive, how many different ways can the six people be divided into three groups so each group shares a car?
  \begin{subanswer}
    % your answer here
  \end{subanswer}

  \item \textbf{Sum with Plus Signs} (10 points)\\
  The digits from 1 to 9 are written in sequence. We can make different sums by adding the "+" symbol between digits. For example, $1+23+456+7+89=576$: The sum is 576, and the greatest number in this case is 456. $12+34+56+78+9=189$: The sum is 189, and the greatest number in this case is 78. Similarly, we can make the sum 90 in various ways. What is the sum of the greatest numbers in these cases?
  \begin{subanswer}
    % your answer here
  \end{subanswer}

  \item \textbf{Functions with Periodicity and Recursion} (10 points)\\
  Find the number of functions $d(x)$ satisfying the following conditions for all $x$:
  \begin{itemize}
    \item $d(x)$ is constant if $n \leq x < n+1$ for any integer $n$;
    \item $d(x-5)=d(x+5)$;
    \item $d(x+1)=d(x)-1$ or $d(x)+1$;
    \item $d(100)=0$.
  \end{itemize}
  \begin{subanswer}
    % your answer here
  \end{subanswer}

  \newpage

  \item \textbf{Shortest Path on a Cone} (10 points)\\
  For the right circular cone below, the radius of the bottom circle is 6, 
  the slant height $AC$ is 24, and $AB=6$. If we draw the shortest path 
  from point $A$ to point $B$ around the cone as shown in the figure, the 
  shortest distance from point $C$ to the path is $k$. What is the value 
  of $10k$?

  \begin{subanswer}
    % your answer here
  \end{subanswer}

  \begin{define}{Combinations with Repetition}
    The equation $x_1+x_2+\cdots+x_m=n$ has $\binom{n+m-1}{n}$ integer solutions such that $x_i \geq 0$ for any $i \in\{1,2, \ldots, m\}$.
  \end{define}



  $$
\binom{n+m-1}{n}
$$

is called a multiset coefficient.
It counts the number of ways to choose $\mathbf{n}$ elements from $\mathbf{m}$ types when repetition is allowed.
Notation sometimes used:

$$
\left(\binom{m}{n}\right)
$$

(read as " $m$ multichoose $n$ ").
So in your equation, the number of solutions equals the number of combinations with repetition.
  
\end{enumerate}

